\section{Matched generalization development on recorded video}
\label{app:generalization-development}
This registered development study investigates the early validation optimum in
the original recorded-video campaign. Three fixed interventions change either
the learning rate (Slow: $3\times10^{-5}$), weight decay (Decay: $0.1$), or the
predictor/context capacity (Compact: width 96, depth 2, context dimensions 16/64).
Each arm trains Framewise, Constant dynamics, ShiftWM (ours), and Action-free at
three seeds for all 30 epochs. The Compact intervention changes these capacity
settings jointly and does not isolate their separate effects.

All inputs, normalization and model selection use the original training and
validation partitions. Neither the original test nor the separately acquired
fresh holdout selects these revisions. Each checkpoint uses the common
all-five-query validation objective. Subsequent five- and ten-block forecasts
are summarized by averaging windows within episodes and then equally averaging
episodes and seeds. Development improvements, if present, require another
independent evaluation for a new confirmatory claim; ordinary optimization and
capacity controls are not presented as additional method novelty.

Table~\ref{tab:generalization-development} reports every matched arm and method.
Tables~\ref{tab:generalization-ci-5} and~\ref{tab:generalization-ci-10} provide
the session/seed paired intervals, including regressions and intervals crossing
zero. Table~\ref{tab:generalization-checkpoints} records the selected epochs,
parameter counts and measured invocation times. All 36 selected packages pass
strict weights-only CPU reload parity on original-validation causal inputs.
Figure~\ref{fig:generalization-summary}a summarizes the six paired endpoint
contrasts alongside the separate horizon-control study in panels b--c.
\input{generated/editorial/validation_controls_landscape_v2_figure.tex}
\begin{table}[p]\centering\footnotesize
\setlength{\tabcolsep}{4pt}\renewcommand{\arraystretch}{1.1}
\begin{tabular}{llrrrr}\toprule
Arm & Method & MSE@5 $\downarrow$ & MSE@10 $\downarrow$ & Gain@5 (\%) & Gain@10 (\%)\\\midrule
Slow & Framewise & 0.15848 & 0.21626 & reference & reference\\
Slow & Constant dynamics & 0.15841 & 0.21605 & \positivegain{+0.04} & \positivegain{+0.09}\\
\rowcolor{orange!9}
Slow & \textbf{ShiftWM (ours)} & 0.15802 & 0.21556 & \positivegain{+0.29} & \positivegain{+0.32}\\
Slow & Action-free & 0.15902 & 0.21564 & -0.34 & \positivegain{+0.28}\\
\midrule
Decay & Framewise & 0.15907 & 0.21689 & reference & reference\\
Decay & Constant dynamics & 0.15899 & 0.21673 & \positivegain{+0.05} & \positivegain{+0.07}\\
\rowcolor{orange!9}
Decay & \textbf{ShiftWM (ours)} & 0.15839 & 0.21948 & \positivegain{+0.43} & -1.20\\
Decay & Action-free & 0.15976 & 0.21742 & -0.44 & -0.25\\
\midrule
Compact & Framewise & 0.15964 & 0.21925 & reference & reference\\
Compact & Constant dynamics & 0.15963 & 0.21923 & \positivegain{+0.01} & \positivegain{+0.01}\\
\rowcolor{orange!9}
Compact & \textbf{ShiftWM (ours)} & 0.15888 & 0.22054 & \positivegain{+0.48} & -0.59\\
Compact & Action-free & 0.16019 & 0.21877 & -0.34 & \positivegain{+0.22}\\
\bottomrule\end{tabular}
\caption{\textbf{Complete matched generalization development study.} All 36 runs completed 30 epochs on original training data. Entries are three-seed, equal-episode validation means; MSE uses fixed train-standardized feature coordinates. Five blocks is the primary development horizon; ten blocks is extrapolation. Slow reduces learning rate, Decay increases weight decay, and Compact jointly reduces predictor/context capacity. Gains compare each method to Framewise within the same arm. \positivegain{Bold green} identifies positive point reductions, not statistical significance. Every arm and method is retained; these are validation results, not an additional held-out test.}
\label{tab:generalization-development}\end{table}

\begin{table}[p]\centering\footnotesize
\setlength{\tabcolsep}{4pt}\renewcommand{\arraystretch}{1.1}
\begin{tabular}{llrrl}\toprule
Arm & Method & $\Delta$MSE & Gain (\%) & Paired 95\% interval\\\midrule
Slow & Constant dynamics & -0.000065 & \positivegain{+0.04} & [-0.000111, -0.000015]\\
\rowcolor{orange!9}
Slow & \textbf{ShiftWM (ours)} & -0.000458 & \positivegain{+0.29} & [-0.000671, -0.000283]\\
Slow & Action-free & +0.000539 & -0.34 & [-0.000062, +0.001162]\\
\midrule
Decay & Constant dynamics & -0.000073 & \positivegain{+0.05} & [-0.000105, -0.000045]\\
\rowcolor{orange!9}
Decay & \textbf{ShiftWM (ours)} & -0.000676 & \positivegain{+0.43} & [-0.001535, -0.000053]\\
Decay & Action-free & +0.000696 & -0.44 & [+0.000208, +0.001192]\\
\midrule
Compact & Constant dynamics & -0.000017 & \positivegain{+0.01} & [-0.000072, +0.000039]\\
\rowcolor{orange!9}
Compact & \textbf{ShiftWM (ours)} & -0.000759 & \positivegain{+0.48} & [-0.001353, -0.000278]\\
Compact & Action-free & +0.000543 & -0.34 & [+0.000029, +0.001149]\\
\bottomrule\end{tabular}
\caption{\textbf{Matched validation uncertainty at 5 blocks.} Differences are method minus Framewise within each arm, in standardized feature MSE; negative favors the named method. Relative gains are $100(\mathrm{Framewise}-\mathrm{method})/\mathrm{Framewise}$. Intervals jointly resample recording sessions and the three training seeds with 10,000 draws. These exploratory development intervals are unadjusted for multiple comparisons; intervals crossing zero do not establish a difference. Positive point gains are colored independently of interval significance.}
\label{tab:generalization-ci-5}\end{table}
\begin{table}[p]\centering\footnotesize
\setlength{\tabcolsep}{4pt}\renewcommand{\arraystretch}{1.1}
\begin{tabular}{llrrl}\toprule
Arm & Method & $\Delta$MSE & Gain (\%) & Paired 95\% interval\\\midrule
Slow & Constant dynamics & -0.000201 & \positivegain{+0.09} & [-0.000322, -0.000086]\\
\rowcolor{orange!9}
Slow & \textbf{ShiftWM (ours)} & -0.000698 & \positivegain{+0.32} & [-0.001144, -0.000339]\\
Slow & Action-free & -0.000610 & \positivegain{+0.28} & [-0.002089, +0.000888]\\
\midrule
Decay & Constant dynamics & -0.000157 & \positivegain{+0.07} & [-0.000241, -0.000079]\\
\rowcolor{orange!9}
Decay & \textbf{ShiftWM (ours)} & +0.002593 & -1.20 & [-0.000748, +0.005568]\\
Decay & Action-free & +0.000532 & -0.25 & [-0.000728, +0.001798]\\
\midrule
Compact & Constant dynamics & -0.000017 & \positivegain{+0.01} & [-0.000194, +0.000153]\\
\rowcolor{orange!9}
Compact & \textbf{ShiftWM (ours)} & +0.001291 & -0.59 & [-0.000945, +0.003425]\\
Compact & Action-free & -0.000483 & \positivegain{+0.22} & [-0.002460, +0.001390]\\
\bottomrule\end{tabular}
\caption{\textbf{Matched validation uncertainty at 10 blocks.} Differences are method minus Framewise within each arm, in standardized feature MSE; negative favors the named method. Relative gains are $100(\mathrm{Framewise}-\mathrm{method})/\mathrm{Framewise}$. Intervals jointly resample recording sessions and the three training seeds with 10,000 draws. These exploratory development intervals are unadjusted for multiple comparisons; intervals crossing zero do not establish a difference. Positive point gains are colored independently of interval significance.}
\label{tab:generalization-ci-10}\end{table}

\begin{table}[p]\centering\footnotesize
\setlength{\tabcolsep}{3.5pt}\renewcommand{\arraystretch}{1.1}
\begin{tabular}{llrrrr}\toprule
Arm & Method & Best epochs & Total (M) & Trainable (M) & Seconds/run\\\midrule
Slow & Framewise & 3/3/3 & 3.763 & 3.469 & 151 [149, 219]\\
Slow & Constant dynamics & 3/3/3 & 3.763 & 3.714 & 159 [159, 234]\\
\rowcolor{orange!9}
Slow & \textbf{ShiftWM (ours)} & 3/3/3 & 3.763 & 3.714 & 161 [161, 236]\\
Slow & Action-free & 3/3/3 & 3.763 & 3.312 & 141 [140, 213]\\
\midrule
Decay & Framewise & 1/1/1 & 3.763 & 3.469 & 148 [146, 218]\\
Decay & Constant dynamics & 1/1/1 & 3.763 & 3.714 & 158 [156, 238]\\
\rowcolor{orange!9}
Decay & \textbf{ShiftWM (ours)} & 2/2/1 & 3.763 & 3.714 & 161 [158, 240]\\
Decay & Action-free & 1/1/1 & 3.763 & 3.312 & 140 [138, 213]\\
\midrule
Compact & Framewise & 2/2/2 & 0.841 & 0.763 & 100 [99, 159]\\
Compact & Constant dynamics & 2/2/2 & 0.841 & 0.828 & 110 [110, 175]\\
\rowcolor{orange!9}
Compact & \textbf{ShiftWM (ours)} & 2/3/3 & 0.841 & 0.828 & 112 [111, 173]\\
Compact & Action-free & 2/1/2 & 0.841 & 0.722 & 93 [93, 149]\\
\bottomrule\end{tabular}
\caption{\textbf{Reusable checkpoints and computation.} Best epochs are ordered by seeds 0/1/2 and selected by the shared all-five-query validation objective after full 30-epoch training. Parameter counts include frozen parameters; trainable counts exclude inactive modules. Seconds/run is the median [minimum, maximum] successful runner invocation, including training, validation and CPU reload verification; prior interrupted work, if any, is excluded. Different GPU partitions are used, so these times are operational evidence, not a hardware-matched speed comparison. All 36 selected packages pass strict CPU reload parity.}
\label{tab:generalization-checkpoints}\end{table}

